\section{Background}
\label{sec:bg}

\subsection{Process Models and Their Silences}

Digital forensic process models describe how an investigation should
proceed, not what classes of evidence exist. The DFRWS framework first
codified the seven-phase pipeline of identification, preservation,
collection, examination, analysis, presentation, and decision~\cite{palmer2001roadmap}.
Carrier and Spafford generalized this into a five-phase \emph{Integrated
Digital Investigation Process} that explicitly aligns the digital phases
with the physical-crime-scene phases that
preceded them~\cite{carrier2003getting}. Beebe and Clark extended the
emphasis on hierarchical objectives, arguing that procedural decomposition
must respond to investigative
objectives rather than artifact taxonomy~\cite{beebe2005hierarchical}.
Carrier's later dissertation~\cite{carrier2006hypothesis} formalized
hypothesis testing as the epistemic core of digital investigation: every
analytical claim is a hypothesis that must be falsifiable from the data.

These models are silent on the question we ask: \emph{what kinds of addressing
does a piece of evidence support, and what are the architectural consequences?}

\subsection{The Three-Source Folk Model and Its Failure Modes}

The implicit, never-formalized folk model of forensic evidence partitions
sources into three: \emph{disk}, \emph{memory}, and
\emph{log}~\cite{kent2006sp80086,casey2011digital,carrier2005file}. Standard
practitioner training and toolchains -- The Sleuth Kit and
Autopsy~\cite{carrier2005file}, Volatility for memory analysis~\cite{ligh2014artmemory,case2017memory},
and timeline tools such as Plaso/log2timeline~\cite{plaso} -- align with this
partition. The folk model has served the field well for two decades but is
breaking under the weight of two developments.

The first is the rise of \emph{enterprise live response} and endpoint
detection-and-response telemetry. When a Velociraptor agent~\cite{velociraptor}
runs a VQL query and returns a result set of currently-open TCP connections,
that result is not on a disk image; it is not yet in any log; it lives,
briefly, only in the network packet that delivers it from the endpoint to
the investigator. Cohen, Bilby, and Garfinkel argued in
2011~\cite{cohen2011distributed} that this kind of capture is a fundamentally
different forensic operation, but the field continued to treat the resulting
data as ``just another log'' -- an analytical convenience that obscures the
property that the data did not exist before the query was issued and, once
captured and transmitted, cannot be modified by an adversary regardless of
subsequent on-host actions.

The second development is \emph{content-addressed} ecosystems. Git,
Sigstore's transparency log~\cite{newman2022sigstore}, in-toto~\cite{torresarias2019intoto},
OCI image manifests~\cite{ociimagespec}, and IPFS~\cite{benet2014ipfs} all
address artifacts by the cryptographic hash of their content; references between
artifacts are themselves hashes. The resulting Merkle DAG has no path, no virtual
address, no timestamp ordering -- only hash $\to$ blob $\to$ content graph.
Supply-chain attacks (SolarWinds~\cite{cisa2021solarwinds},
\texttt{xz-utils}~\cite{cve2024xz}) forced the forensic community to reckon
with evidence trails that live entirely in content-addressed stores that no
traditional triage tool addresses.

\subsection{The Gap}

The literature provides three separate strands -- traditional disk and log
forensics, memory forensics as it matured through Volatility and
Rekall~\cite{ligh2014artmemory,case2017memory}, and emerging work in live
response and supply-chain forensics -- that have not been unified by a
common formal model.

The closest representation-level effort is the CASE/UCO
ontology~\cite{caseuco2018,casey2017advancing}, which provides a typed
vocabulary for describing digital evidence \emph{after} it has been parsed.
Our taxonomy is complementary, not competing: CASE/UCO is \emph{annotation}
(the schema in which findings are exchanged), while the navigation-primitive
taxonomy is \emph{architecture} (the addressing discipline by which findings
are produced). A CASE/UCO record consumed by a downstream tool conveys the
\emph{what}; a navigation-primitive tag (\eg{}, \texttt{nav-primitive: Q})
conveys the \emph{how}, which CASE/UCO does not encode and which determines
the trust profile and durability claims that the consumer can make.

The closest infrastructure-level effort is Carrier's data-abstraction model
and its concrete instantiation in The Sleuth Kit (TSK)~\cite{carrier2005file,tsk_repo}.
TSK exposes \texttt{img\_*}, \texttt{vol\_*}, \texttt{fs\_*}, and \texttt{file\_*}
APIs corresponding to a layered traversal: image $\to$ volume $\to$
filesystem $\to$ file. This layering is foundational for \src{P} navigation
specifically: TSK's \texttt{img\_} layer is what we generalize as CONTAINER,
and TSK's \texttt{fs\_}/\texttt{file\_} layers together constitute the
\src{P}-specialization of NAVIGATE in our framework. The crucial difference
is that TSK's layering is \emph{specific} to disk evidence -- it has no
analog of \texttt{img\_} for memory dumps, no analog of \texttt{vol\_} for
log streams, no analog of \texttt{file\_} for hash-addressed blobs. Our
taxonomy generalizes Carrier's layered insight to four other addressing
schemes that did not exist in their current form when TSK was designed.

Garfinkel's 2009 work on automating disk forensic processing with TSK, XML,
and Python~\cite{garfinkel2009automating} provides early evidence that even
within \src{P} alone the absence of a uniform addressing model forces
ad-hoc adapters: Garfinkel's DFXML output exists precisely because tool-output
schemas were not interoperable. The same problem at five-source scale
motivates the present paper.

Garfinkel's 2010 agenda~\cite{garfinkel2010digital} identified cross-source
correlation as one of the key challenges of the coming decade; sixteen
years later, the field has rich tooling for each source individually but no
formal account of what ``cross-source'' means. We aim to provide that
account.
